\section{Related Work}

\textbf{Active Feature Acquisition}\quad Prior works in AFA \citep{saar2009active, sheng2006feature} study the trade-off between prediction and feature acquisition cost in a non-longitudinal context. Recent works \citep{shim2018joint, yin2020reinforcement, janisch2020classification} frame AFA as an RL problem; however, they are challenging to train in practice due to the complicated state and action space with dynamically evolving dimensions. Other works \citep{li2021active} use surrogate generative models, which model multidimensional distributions, or employ greedy policies that ignore joint informativeness (some features are only informative in combination) \citep{covert2023learning, ma2018eddi, gong2019icebreaker}. To overcome these issues, \citet{valancius2024acquisition} employ a non-parametric oracle-based method, and  \citet{ghosh2023difa} learn a differentiable acquisition policy by jointly optimizing selection and prediction models through end-to-end training. However, these works address AFA in a non-temporal context, whereas many real-world tasks require sequential decisions over time, motivating the need for \emph{longitudinal} AFA.

% While these methods focus on single-timepoint acquisition, many real-world applications require decisionsto be made sequentially over time, motivating the longitudinal AFA setting.
%Alternatively, \citep{gadgil2023estimating} considers a discriminative approach that directly estimates conditional mutual information between the estimated information change and the true information differences. 


\textbf{Longitudinal Active Feature Acquisition}\quad Longitudinal AFA extends the standard AFA to include the temporal dimension. \quad
% Complementary work in the signal‑processing community studies active state tracking, optimizing sensing policies under explicit cost constraints; e.g., \citep{zois2017active} derives dynamic‑programming and low‑complexity myopic policies for finite‑state Markov chains. 
In this setting, a policy must account for temporal constraints where past time points can no longer be accessed. As a result, a policy must decide which features to acquire and determine the optimal timing for each acquisition. For instance, \citet{zois2017active} derive dynamic‑programming and low‑complexity myopic policies for finite‑state Markov chains, ASAC \citep{yoon2019asac} uses an actor-critic approach to jointly train a feature selection and a prediction network. Additionally, \citet{nguyen2024active} use RL to optimize acquisition timing, but acquire all available features at the selected time point. Recently, \citet{qin2024risk} prioritize timely detection by allowing flexible follow-up intervals in continuous-time settings, ensuring that predictions are timely and accurate. However, RL-based approaches can be challenging to train due to issues such as sparse rewards, high-dimensional observation spaces, and complex temporal dependencies \citep{li2021active}. Rather than relying on RL to learn a policy, our proposed approach, \texttt{\OurMethod}, explicitly 
plans by scoring future candidate acquisitions using either \texttt{NOCT-Contrastive} or \texttt{NOCT-Amortized}.\looseness-1

